% Generated draft: requires mathematical and source review.
\documentclass[11pt]{article}
\usepackage[margin=25mm]{geometry}
\usepackage{fontspec}
\setmainfont{Latin Modern Roman}
\usepackage{amsmath,amssymb,mathtools,unicode-math}
\setmathfont{Latin Modern Math}
\usepackage{enumitem,needspace,xurl,hyperref}
\setlength{\parindent}{0pt}
\setlength{\parskip}{5pt}
\setlength{\emergencystretch}{3em}
\newcommand{\N}{\mathbb{N}}
\newcommand{\Z}{\mathbb{Z}}
\newcommand{\Q}{\mathbb{Q}}
\newcommand{\R}{\mathbb{R}}
\newcommand{\C}{\mathbb{C}}
\newcommand{\F}{\mathbb{F}}
\newcommand{\A}{\mathbb{A}}
\newcommand{\PP}{\mathbb{P}}
\newcommand{\E}{\mathbb{E}}
\newcommand{\eps}{\varepsilon}
\newcommand{\dd}{\,\mathrm{d}}
\begin{document}
\section*{UnsolvedMath ID: 30003719}
\textbf{Source problem code:} catalogue entry
\subsection{Definitions and mathematical statement}
Principal component analysis (PCA) seeks lower-dimensional descriptions of data. Here the intended data space is a manifold, and the proposed analysis uses covariance tensor fields. The supplied record does not specify the manifold's geometric structure, the data distribution or sampling model, or the definition and tensor type of the covariance field. Consequently, these terms do not yet determine a particular mathematical construction.

The source poses the following related research questions:
\begin{enumerate}
\item Formulate a backward nested procedure for PCA on manifolds using covariance tensor fields. Specify the admissible component models, the meaning of nesting, and the rule for selecting successive models.
\item Characterize when that backward procedure agrees with the corresponding forward procedure, termed \emph{commutative PCA} in the supplied record. Agreement requires a specified forward procedure and a specified comparison of its outputs with the backward outputs, including treatment of nonunique choices; these are not supplied.
\item Determine conditions under which PCA or classification is appropriate when curvature is substantial. The relevant curvature quantity, the scale at which curvature is substantial, the classification setting, and the criterion of appropriateness are unspecified.
\item Determine whether an analysis across multiple scales, analogous to that available in the linear setting, can be developed for these manifold procedures. The scale parameter, the intended linear comparison, and the properties required of such a theory are unspecified.
\end{enumerate}
These questions constitute a research programme rather than a fully specified proposition. No universal quantifier over manifolds, covariance fields, data laws, or classification tasks can be extracted from the supplied record without additional hypotheses. In particular, no particular metric, loss function, projection, component family, or notion of scale is assumed in this draft.
\subsection{Short English statement}
Develop backward nested PCA from covariance tensor fields on manifolds, characterize agreement with forward PCA, and determine suitable uses of PCA or classification in substantial curvature and the possibility of a multiscale theory analogous to linear PCA. The supplied record leaves the geometric, statistical, and algorithmic definitions unresolved, so it does not determine a unique mathematical proposition.
\subsection{Sources}
\begin{itemize}
\item Local UnsolvedMath problem registry, permanent ID \texttt{30003719}, title \emph{Nested Principal Components from Manifold Covariance Fields}, exact source problem code \texttt{catalogue entry}. The record attributes the questions to Oberwolfach Reports workshop discussions, records proposed year 2018, and supplies no individual proposer.
\item Supplied report citation: \emph{Statistics for Data with Geometric Structure} (2019), Oberwolfach Reports, pp.\ 125--186, DOI \texttt{10.4171/owr/2018/3}. \url{https://doi.org/10.4171/owr/2018/3}. The background describes Oberwolfach Reports as documenting mathematical workshops and their open-problem discussions. The citation's 2019 year and the record's 2018 report identifiers are retained without reconciliation.
\item Supplied literature source: \emph{Oberwolfach Report 3/2018}, \url{https://doi.org/10.4171/owr/2018/3}. Literature checked at: \texttt{2026-08-21}. This is supplied dated curation, not an independently verified assessment. It labels the problem open and states that existing special-case constructions do not settle nested or commuting manifold PCA, curvature-robust classification, or a general multiscale analysis. The supplied background expressly calls for independent verification of time-sensitive status and cited claims.
\item The registry labels the statement \texttt{corrected\_verified}, explaining that mathematical content was preserved while notation, typography, and sentence structure were repaired for publication. This is supplied editorial metadata, not independent mathematical verification of this draft.
\end{itemize}
\end{document}
